\chapter{Command-line interface}\label{app:cli}

\NF{} installs one console command, \code{nurbforge-llm}, an interactive shell for building
geometry from natural language with a local model served by Ollama (\cref{ch:llm}). It is a
small wrapper around \apiref{llm.GeometryAgent}. Everything it does is also available from
Python.

\section{Synopsis}

\begin{lstlisting}[style=shell]
nurbforge-llm [-h] [--model MODEL] [--host HOST] [--out OUT]
              [--mesh-size MESH_SIZE] [--verify-rounds VERIFY_ROUNDS]
\end{lstlisting}

The command is equivalent to \code{python -m nurbforge.llm.cli}.

\section{Options}

\begin{longtable}{@{}p{0.25\textwidth}p{0.30\textwidth}p{0.39\textwidth}@{}}
\toprule
\textbf{Option} & \textbf{Default} & \textbf{Meaning} \\
\midrule
\endhead
\code{-h}, \code{--help} & & Print the usage and exit. \\
\code{--model MODEL} & \code{NURBFORGE_LLM_MODEL}, else \code{qwen2.5:14b} & Name of the Ollama
  model. It must support tool calling and be pulled (\code{ollama pull MODEL}). \\
\code{--host HOST} & \code{OLLAMA_HOST}, else \code{http://localhost:11434} & URL of the Ollama
  server. A host without scheme (\code{127.0.0.1:11434}) gets \code{http://} prepended. \\
\code{--out OUT} & none & On exit, mesh all \emph{surfaces} (\code{NurbsSurface} objects) in
  the workspace with gmsh and write them to this file. The format follows from the extension
  (\code{.msh}, \code{.vtk}, \code{.stl}, ...). Needs gmsh. Shells and curves are not
  meshed. Export shells from the session with the \toolref{export_step} tool instead. \\
\code{--mesh-size MESH_SIZE} & gmsh default & Global maximum element size for \code{--out}, in
  model units. \\
\code{--verify-rounds N} & 1 & Verification rounds after each answer. Every object created or
  changed by the request is measured and the report is sent back to the model, which repairs
  or confirms the result. 0 disables verification. \\
\bottomrule
\end{longtable}

\section{Environment variables}

\begin{tabularx}{\textwidth}{@{}lX@{}}
\toprule
\textbf{Variable} & \textbf{Effect} \\
\midrule
\code{OLLAMA_HOST} & Default server address for \code{nurbforge-llm} and for
  \apiref{llm.OllamaClient} created without a \code{host} argument. It is read when
  \code{nurbforge.llm} is imported. \\
\code{NURBFORGE_LLM_MODEL} & Default model name, read the same way. \\
\bottomrule
\end{tabularx}

\section{Interactive use}

After start-up the shell prints a banner with model and host, then reads one request per line
after the prompt \code{>}:

\begin{itemize}
  \item Each request is sent to the agent. The model may call any number of tools (up to 20 steps
    per request), then answers in text. The answer is printed, followed by the tool calls of
    the request, one per line, as \code{[tool] arguments -> result}.
  \item The workspace persists across requests. Objects can be referred to by name in later
    requests (``now mirror \code{plate} at the yz plane'').
  \item An empty line or end of input (Ctrl-D, on Windows Ctrl-Z Enter) ends the session. If
    \code{--out} is given, the surfaces are meshed and written, and the number of surfaces is
    reported.
\end{itemize}

\section{Examples}

\begin{lstlisting}[style=shell]
# default model and server
nurbforge-llm

# a different model on another machine, mesh the surfaces on exit
nurbforge-llm --model llama3.1:8b --host http://gpu-server:11434 --out model.msh --mesh-size 0.5

# without verification rounds (faster, less reliable)
nurbforge-llm --verify-rounds 0

# configuration through the environment
export OLLAMA_HOST=http://gpu-server:11434
export NURBFORGE_LLM_MODEL=qwen2.5:32b
nurbforge-llm
\end{lstlisting}

A session might look like this (the model's wording varies):

\begin{lstlisting}[style=text]
Nurbforge LLM shell (qwen2.5:14b @ http://localhost:11434). Empty line to quit.
> Create a box named plate from (0,0,0) with size 40 x 30 x 5 and drill a hole of radius 4
  through its centre along z.
...
  [solid] {'name': 'plate', 'kind': 'box', ...} -> stored plate: ...
  [solid] {'name': 'drill', 'kind': 'cylinder', ...} -> stored drill: ...
  [boolean] {'name': 'plate', 'a': 'plate', 'b': 'drill', 'op': 'difference'} -> ...
> Export plate to plate.step
...
>
\end{lstlisting}

The tools available to the model, with their parameters, are listed in \cref{app:llmtools}.
